\documentclass[11pt,a4paper]{article}
\usepackage[T1]{fontenc}
\usepackage[utf8]{inputenc}
\usepackage{amsmath,amssymb,amsthm}
\usepackage[margin=1in]{geometry}
\usepackage[colorlinks=true,linkcolor=blue,urlcolor=blue]{hyperref}
\theoremstyle{plain}
\newtheorem{theorem}{Theorem}
\newtheorem{lemma}[theorem]{Lemma}
\newtheorem{proposition}[theorem]{Proposition}
\newtheorem{corollary}[theorem]{Corollary}
\theoremstyle{definition}
\newtheorem{remark}[theorem]{Remark}

\title{Recovering the unwritten recurrences of the table OEIS A236789}
\author{Adrian Perez Fontelles\\ \small Independent researcher}
\date{2 September 2026}
\begin{document}
\maketitle

\begin{abstract}
OEIS A236789 is a two-dimensional table $T(n,k)$ whose empirical recurrences are, for several of
its lines, given only as an ORDER --- ``k=2: [order 35]'' and the like --- with the
coefficients in a linked file rather than in the entry. Those conjectures can still be settled,
and the recurrences recovered. A line of the table is a fixed-width or fixed-height array
count, hence a walk count on finitely many vertices, hence satisfies some monic recurrence of
order at most that number; Berlekamp--Massey on twice as many exact terms returns the MINIMAL
one. Where its order equals the order the entry states --- and where the same holds after the
first few terms are dropped --- any recurrence of that order the line satisfies has a
characteristic polynomial that is a multiple of the minimal one and of equal degree, hence is
that polynomial. This note settles column 2, column 3.
\end{abstract}

\noindent\small 2020 Mathematics Subject Classification. 05A15, 11B37, 15A18.\normalsize

\section{The table and the unwritten conjectures}

OEIS A236789 is ``T(n,k)=Number of (n+1)X(k+1) 0..2 arrays with the upper median plus the lower median minus the minimum of every 2X2 subblock equal.''. It is read by antidiagonals and begins
\[
81, 255, 255, 770, 1097, 770, 2475, 4452,\ \dots
\]
The entry carries, among its empirical claims:
\begin{quote}\small
k=2: [order 35]\\
k=3: [order 83]
\end{quote}
As of the ``Last modified'' line on the live entry (Jul 23 2025 09:33:04, revision 6) these are still
recorded as empirical, and the recurrences themselves are not in the entry text.

\section{Why a line of the table is a walk count}

Fix the index that the line fixes. The entry defines $T(n,k)$ as a number of arrays of a shape
determined by $n$ and $k$, subject to a condition local to a bounded window of consecutive
lines. Substituting the fixed index into the entry's own wording turns the definition into an
ordinary one-parameter array count, and for such a count the admissible windows are the
vertices of a finite digraph and an array is a walk. So the line is a walk count on $S$
vertices, and by the Cayley--Hamilton theorem it satisfies a monic integer recurrence of order
at most $S$.

\section{Recovering the recurrences}

\begin{lemma}\label{lem:bm}
A sequence known to satisfy some linear recurrence of order at most $S$ is determined, with its
minimal recurrence, by its first $2S$ terms; Berlekamp--Massey applied to them returns that
minimal recurrence exactly.
\end{lemma}

\subsection*{Column $k=2$}
Substituting the fixed index into the entry's wording gives an ordinary one-parameter array
count, a walk on $S=82$ vertices. Berlekamp--Massey on $2S=164$ exact terms
returns a minimal recurrence of order $35$ --- the order the entry states --- with
integer coefficients
\[
(c_1,\dots,c_{35})=(22,\ -145,\ -234,\ 6880,\ -19940,\ -92792,\ 593404,\ 66753,\ -7105820,\ 10259982,\ 44201366,\ -120008970,\ -136396696,\ 699300702,\ 46313296,\ -2478218108,\ 1244671494,\ 5661814333,\ -5050034688,\ -8465863134,\ 10417098396,\ 8197726584,\ -13227353752,\ -4935036828,\ 10838034848,\ 1645800008,\ -5758739808,\ -170986944,\ 1942394944,\ -64872256,\ -395392512,\ 22059776,\ 43732480,\ -1955840,\ -1996800).
\]
so that $a(m)=\sum_{i=1}^{35}c_i\,a(m-i)$ for every $m$. Removing the first
$35+4$ terms and repeating gives order $35$ again, which is what the
identification below needs.

\subsection*{Column $k=3$}
Substituting the fixed index into the entry's wording gives an ordinary one-parameter array
count, a walk on $S=207$ vertices. Berlekamp--Massey on $2S=414$ exact terms
returns a minimal recurrence of order $83$ --- the order the entry states --- with
integer coefficients
\[
(c_1,\dots,c_{83})=(28,\ -138,\ -3131,\ 35138,\ 98877,\ -2800396,\ 3228326,\ 125292435,\ -409487545,\ -3658343046,\ 18570209996,\ 73890433500,\ -537953433565,\ -1036914148304,\ 11353606846264,\ 9262834011851,\ -185372213922877,\ -25208328378633,\ 2425698908549131,\ -730637245299426,\ -26048490590020008,\ 15143385673379179,\ 233451188720055518,\ -174825791573522016,\ -1767652157336689202,\ 1467431964468868485,\ 11407789258412574883,\ -9623770890253871944,\ -63127434706444524945,\ 50770128670992361147,\ 300650993083797005380,\ -218117763864362502635,\ -1234668756586394898319,\ 764888665201445819942,\ 4373980839183196365696,\ -2173683372325258383932,\ -13359164252517494081969,\ 4896179587664924611430,\ 35129294952621730751682,\ -8242478561670752859501,\ -79382846980586487706211,\ 8389524523571666152049,\ 153797473431616303838547,\ 2755242749977941205912,\ -254779186886689729303186,\ -34481742380708237124529,\ 359745868383003112593830,\ 88890095748878826722233,\ -431313002884093671200740,\ -153049379398767957994553,\ 437009949663734140734854,\ 201687752489851679515320,\ -371871546416467790083348,\ -212599809044824918936599,\ 263490591654515798611976,\ 182559643848927772206706,\ -153490862386214424053150,\ -128660670405546124340154,\ 72006378868591688783024,\ 74532873527896704989144,\ -26172691324997580887576,\ -35393400947227810195624,\ 6719942631127418614072,\ 13685459502332757231440,\ -821383510292776426168,\ -4260328525115035551824,\ -209934595568546184240,\ 1048937671203428433024,\ 154128775566146657024,\ -198531253345325931008,\ -48204448855190652608,\ 27487303581507030144,\ 9727879220036788864,\ -2505044070247747584,\ -1330974243904608768,\ 103215506373133312,\ 118974579224262656,\ 5363680562884608,\ -6098940012625920,\ -894945553465344,\ 110765149716480,\ 34345987080192,\ 2192984506368).
\]
so that $a(m)=\sum_{i=1}^{83}c_i\,a(m-i)$ for every $m$. Removing the first
$83+4$ terms and repeating gives order $83$ again, which is what the
identification below needs.

\begin{theorem}
For each line above, the displayed recurrence holds for every index, and it is the recurrence
the entry records.
\end{theorem}

\begin{proof}
It holds by Lemma~\ref{lem:bm}. For the identification, the entry asserts that the line
satisfies SOME recurrence of the stated order, possibly only from some index on. Let $q$ be its
characteristic polynomial and $q_{\min}$ the minimal polynomial of the tail. Then
$q_{\min}\mid q$, and the second Berlekamp--Massey run --- on the line with its first few
terms removed --- shows $\deg q_{\min}$ equals the stated order, which is $\deg q$. Both are
monic, so $q=q_{\min}$.
\end{proof}

\section{Verification}

Three checks.

First, each line's model was compared against the entry's own published values for that line,
taken from the antidiagonal DATA read in either orientation and from the ``Table starts''
block, and agrees at every available term. This is the only point at which reading the entry's
English enters, and a misreading gives different counts.

Second, each recovered recurrence was evaluated directly on the model's values, with no
Berlekamp--Massey involved, and holds at every index tested.

Third, each order was computed twice by independent routes --- modulo a $61$-bit prime, where
every intermediate is a machine word, and again in exact rational arithmetic --- and once more
on the line with its first few terms dropped, which is what the identification requires. All
agree.

\begin{thebibliography}{9}
\bibitem{oeis} The OEIS Foundation, \emph{The On-Line Encyclopedia of Integer
Sequences}, \url{https://oeis.org}, 2026. Sequence A236789.
\bibitem{massey} J.~L.~Massey, Shift-register synthesis and BCH decoding, \emph{IEEE
Trans. Inform. Theory} \textbf{15} (1969), 122--127.
\bibitem{stanley} R.~P.~Stanley, \emph{Enumerative Combinatorics}, Volume 1, 2nd ed.,
Cambridge University Press, 2011. (Transfer-matrix method, Section 4.7.)
\bibitem{horn} R.~A.~Horn and C.~R.~Johnson, \emph{Matrix Analysis}, 2nd ed., Cambridge
University Press, 2013.
\end{thebibliography}

\end{document}
